\chapter{Mechanizing Floating Point Bitblasting}
\label{ch:floating-point}

\paragraph*{Attribution.} This chapter is based on \emph{Mechanizing Floating
Point Bitblasting}~\citep{bhat2027floatingpoint}, currently in submission.

\lettrine{F}{loating-point} verification is central to trustworthy numerical
computing. Existing approaches fall into two camps: proof assistants that
provide precise formalizations of floating-point semantics, such as
Flocq~\citep{boldo2011flocq}, but require manual proofs; and SMT solvers such
as cvc5~\citep{barbosa2022cvc5} and Bitwuzla~\citep{niemetz2023bitwuzla} that
perform efficient, push-button reductions of floating point to bitvectors via
bitblasting with a ``soft floating-point unit.'' The latter approach relies
on the remarkable progress in bitvector SMT solving---including the verified
bitblasting infrastructure of \cref{ch:background}---but the floating-point
reductions themselves are subtle, complex, and had never been verified.

This chapter bridges that gap by presenting the first \emph{verified}
floating-point bitblaster: a library and associated tactic in the Lean proof
assistant that
(a)~mechanizes the SMT-LIB theory of floating point (\qffp{}),
(b)~extends Lean's bitvector solver \bvdecide{} with floating-point
bitblasting capabilities,
(c)~proves the correctness of the bitblasting reduction relative to the
SMT-LIB theory, and
(d)~carefully builds the underlying circuits to support the small
floating-point formats (e.g., E4M3 and E5M2) used in machine-learning
quantization.

Our evaluation shows that verified bitblasting is practical: on the SMT-COMP
benchmarks our solver runs within a small geometric-mean slowdown of
Bitwuzla, while offering end-to-end correctness guarantees. On new benchmarks
targeting small floating-point formats, we are substantially faster than the
sound exhaustive-enumeration baseline, establishing the first sound, verified
solver for these formats.

\section{Algorithmic Overview: How to Bitblast a Floating Point Multiply}
\label{sec:fp:overview}

\sid{Port from papers/paper-floating-point-in-lean/paper.tex \S 2 (lines
577--774). Figures: chapters/floating-point/images/ (overview.pdf,
floatToFixed.pdf, trimming.pdf, sqrt\_pitfalls.pdf).}

\section{A Full Mechanization of the SMT-LIB QF\_FP Theory}
\label{sec:fp:spec}

\sid{Port from papers/paper-floating-point-in-lean/paper.tex \S 3 (lines
775--1037).}

\section{Verified Floating Point Bitblasting Circuits for QF\_FP Operations}
\label{sec:fp:design}

\sid{Port from papers/paper-floating-point-in-lean/paper.tex \S 4 (lines
1038--1386).}

\section{From Verified Circuits to a Solver}
\label{sec:fp:solver}

\sid{Port from papers/paper-floating-point-in-lean/paper.tex \S 5 (lines
1387--1604).}

\section{Exhaustive and Randomized Testing for Specification Fidelity}
\label{sec:fp:testing}

\sid{Port from papers/paper-floating-point-in-lean/paper.tex \S 6 (lines
1605--1750).}

\section{Case Study: Gaps Discovered by Mechanization}
\label{sec:fp:surfaced}

\sid{Port from papers/paper-floating-point-in-lean/paper.tex \S 7 (lines
1751--1790).}

\section{Quantitative Evaluation of Solver Performance}
\label{sec:fp:evaluation}

\sid{Port from papers/paper-floating-point-in-lean/paper.tex \S 8 (lines
1791--2002). Plots: chapters/floating-point/plots/smtlib-rand/ and
chapters/floating-point/plots/instcombine-small/ (smtlib-cactus.pdf +
smtlib-stats.tex). NB: the paper's \textbackslash inputs of
plots/smtlib-rand/outputs/\{triple,config\} do not exist in the repo; skip
them.}

\section{Related Work}
\label{sec:fp:related}

\sid{Port from papers/paper-floating-point-in-lean/paper.tex \S 9 (lines
2064--2182).}

\section{Conclusions, Limitations, and Future Work}
\label{sec:fp:conclusion}

\sid{Port from papers/paper-floating-point-in-lean/paper.tex \S 10 (lines
2183--).}
